\section{From allocation to silence}
\label{sec:population}
\subsection{Arrival exposure and the uniformity requirement}
An allocation rule is a statement about one decision process. A Fermi
explanation requires a distribution of such processes, their launch
histories, and propagation to the region where absence is being explained.
Let $N_i^{\rm act}(t)$ denote outward-capable autonomous nodes and
$\lambda_i(t)$ their initiation rates per node-time. An independent
Poisson approximation gives a surviving-front mean
\begin{equation}
 \mu=\sum_i\int\lambda_i(t)N_i^{\rm act}(t)[1-q_{R,i}(t)]\dd t.
 \label{eq:mu}
\end{equation}
Here $q_R$ is an extinction probability. For a Galton--Watson process it
is the smallest fixed point of the offspring generating function; for
Poisson offspring of mean $\Rogue$,
$q_R=\exp[\Rogue(q_R-1)]$ \citep{Harris1963}. Except for degenerate
immortal one-descendant processes, a mean at most one gives extinction
almost surely. Survival is not arrival, and extinct fronts can still
leave finite artefacts. Replace $\mu$ by $\mu_A(z,T)$ when launch,
survival, and propagation into the target region by time $T$ have all
been accounted for.

For a fixed independent list of civilizations with Poisson arrival
means $\mu_{A,i}$, $P_0^{\rm fixed}=\exp(-\sum_i\mu_{A,i})$.
If origins themselves form a Poisson process with intensity measure
$\nu(\dd z)$, independently marked by their arrival means, the void
probability instead is
\begin{equation}
 P_0^{\rm origins}
 =\exp\left[-\int(1-\exp[-\mu_A(z,T)])\nu(\dd z)\right].
 \label{eq:origins}
\end{equation}
A random number of parent civilizations clusters fronts and changes the
likelihood. For two fixed civilizations with mean two fronts each the
probability is $0.0183$; for a Poisson number of mean two it is $0.1774$.
These archived numbers are retained. The derivation and a new simulation
are given in \cref{app:population-proof}.

\begin{proposition}[Necessary uniformity and the exact two-type window]
\label{prop:uniformity}
Let $0<\bar N_A=\int\nu(\dd z)<\infty$, $0<P_*<1$, and
$\mathcal L_*=-\ln P_*$. Suppose a fraction $\pi_{\rm out}$ of origins under the
normalized arrival-relevant intensity has $\mu_A\geq\mu_{\rm out}>0$.
Then $P_0\geq P_*$ requires
\begin{equation}
 \pi_{\rm out}\leq
 \frac{\mathcal L_*}{\bar N_A[1-\exp(-\mu_{\rm out})]}.
 \label{eq:uniformity-bound}
\end{equation}
For the complementary class labelled inward, let $c_{\rm in}=\mathbb E[1-\exp(-\mu_A)\mid\text{inward}]$.
Set $c_{\rm in}=0$ when that class has zero probability. Then also $(1-\pi_{\rm out})c_{\rm in}\leq \mathcal L_*/\bar N_A$.
For two types with exactly specified means
$\mu_{\rm out}>\mu_{\rm in}\geq0$, write
$f_k=1-\exp(-\mu_k)$. The exact necessary and sufficient condition is
\begin{equation}
 f_{\rm in}+\pi_{\rm out}(f_{\rm out}-f_{\rm in})
 \leq \mathcal L_*/\bar N_A.
 \label{eq:two-type-window}
\end{equation}
Without any upper bound on the means, $\bar N_A\leq \mathcal L_*$ is sufficient
for $P_0\geq P_*$. The weaker condition
$\bar N_Af_{\rm out}\leq \mathcal L_*$ is sufficient for arbitrary mixtures only
in the exactly specified two-type model, not when $\mu_{\rm out}$ is
merely a lower bound.
\end{proposition}
\noindent\emph{Proof sketch.}
The negative logarithm of \cref{eq:origins} is
$\bar N_A\mathbb E[1-\exp(-\mu_A)]$. Bound its non-negative outward
and inward contributions separately; evaluate the expectation exactly
for two types. A complete proof and boundary cases are in
\cref{app:population-proof}.

At $P_*=0.05$ and $\mu_{\rm out}=3$, the necessary upper outward
fractions for $\bar N_A=10^2,10^3,10^4$ are respectively
$3.1527\%$, $0.31527\%$, and $0.031527\%$. These are conditional
compatibility bounds, not posterior probabilities. For example, a prior
on the number of origins could make large exposure unlikely; then
uniform inward behaviour need not be invoked. Conversely, with a large
arrival-relevant population, a small expansionist tail can exhaust the
entire permitted exposure. Any inward leakage makes the constraint on
that tail tighter, not looser.

The archived design is far from a near-uniform all-inward allocation:
its reference count is $219/1024$, and its design-measure sensitivity is
reported in the supplement. But $805$ failed all-inward tests cannot be
identified with $805$ replicating civilizations. Reoptimizing observation
on the no-replication face, with exactly the same frozen inputs and no
new dormant-backup technology, gives $421/1024$ no-replication optima.
Of these, $202$ use observation without replication; $603/1024$ prefer
replication. These additive calculations strengthen the logical
distinction rather than providing an astronomical prevalence estimate.
Even this selected synthetic ensemble is nowhere near the high-exposure
uniformity requirement under a literal population interpretation, which
its sampling design does not justify.

Under a uniform external price multiplier $\kappa$, the archived statistic
$\rho$ gives $\Pr(\rho>\kappa)$ for loss of the \emph{all-inward} corner.
It equals an outward population fraction only after specifying a
population measure and identifying every such failure with the relevant
outward behaviour. With observation or backups free, the face optimizer
must be recalculated; the old $\rho$ is not the no-replication threshold.
This is why a change in wording alone cannot carry the old sensitivity
experiment into the new population claim.

\subsection{Preemption: a game rather than an adjusted corner}
To formalize strategic complementarity, let $N\geq2$ civilizations choose
$x_i\in[0,\bar X]$. Suppose
\begin{equation}
 U_i=f(x_i)+\iota x_i\sum_{k\ne i}x_k,
 \qquad \iota\geq0,
 \label{eq:preemption-game}
\end{equation}
where $f$ is continuously differentiable and strictly concave. It can
include the static local-versus-external objective, including the base
residual slope $S_{j0}$. This is a specified payoff externality, not a
claim that fear of neighbours has a universal linear form. Competition
among expanding lineages is an established motivation for such a test
\citep{Hanson1998}.

\begin{proposition}[Preemption equilibria and a separate cascade test]
\label{prop:preemption}
In \cref{eq:preemption-game}, a profile is a pure Nash equilibrium exactly
when, with $X_{-i}=\sum_{k\ne i}x_k$,
\begin{equation}
 \begin{cases}
 f'(0)+\iota X_{-i}\leq0,&x_i=0,\\
 f'(x_i)+\iota X_{-i}=0,&0<x_i<\bar X,\\
 f'(\bar X)+\iota X_{-i}\geq0,&x_i=\bar X.
 \end{cases}
 \label{eq:preemption-KKT}
\end{equation}
A pure equilibrium exists. The all-inward profile is an equilibrium iff
$f'(0)\leq0$; the all-cap profile is an equilibrium iff
$f'(\bar X)+\iota(N-1)\bar X\geq0$. When both hold there are at least
two equilibria. Symmetric interior equilibria are exactly the roots of
$f'(x)+\iota(N-1)x=0$ in $(0,\bar X)$.

For the distinct binary-action game $x_i\in\{0,\bar X\}$, put
$d_{\rm bin}=[f(\bar X)-f(0)]/\bar X$. Its all-inward and all-outward
profiles are equilibria iff, respectively,
$d_{\rm bin}\leq0$ and $d_{\rm bin}+\iota(N-1)\bar X\geq0$.
With inward tie-breaking, one externally maintained outward player
triggers all other players under asynchronous best-response updates
iff $d_{\rm bin}+\iota\bar X>0$.
\end{proposition}
\noindent\emph{Proof sketch.}
Strict concavity gives the endpoint and interior derivative conditions.
The potential $\sum_i f(x_i)+\iota\sum_{i<k}x_ix_k$ has a maximum on
the compact strategy cube. Binary choices instead compare finite payoff
differences; they cannot be tested using $f'(0)$ alone. Full proofs are
in \cref{app:preemption-proof}.

Failure of the inward derivative test at maximal rival investment does
not establish an outward equilibrium. For $N=2$, $\bar X=1$,
$f(x)=-x-5x^2$, and $\iota=2$, that derivative is $-1+2=1>0$,
but the only equilibrium is $(0,0)$. Nor do two stable endpoint
equilibria imply a one-defector cascade. In the binary game with five
players, $d_{\rm bin}=-2.5$ and $\iota=\bar X=1$, both endpoints are
stable, but a single defector gives the remaining players incentive
$-1.5<0$. A sufficiently large cluster can matter when one player does
not. Population uniformity is therefore an equilibrium-selection and
containment problem as well as a distribution of individual preferences.

Virtual population growth affects that problem only through physical
agency. If worlds grow as $\exp(g_Wt)$, their outward-capable fraction
falls as $\exp(-g_ct)$, and the launch-times-arrival factor stays bounded
above and below by positive constants, integrated infinite-horizon
exposure is finite exactly when $g_c>g_W$. This is the elementary
integrability criterion for $\int_0^\infty\exp[(g_W-g_c)t]\dd t$.
Without actuator access the exposure can instead be zero, regardless
of the virtual population. Recursive population count alone is neither
a launch rate nor a guarantee of containment.

\subsection{Visibility is a separate inverse problem}
Neither an allocation share nor an arrival probability fixes spectrum,
power, geometry, or duty cycle. With mutually exclusive source classes,
a schematic Poisson detection model is
\begin{equation}
 \Lambda_j(T)=\int_0^T R_{\rm civ}(t_e)
 \mathcal P_j(T-t_e\mid t_e)\pi_j(T,t_e)\dd t_e,
 \qquad P(0)=\exp[-\sum_j\Lambda_j(T)].
 \label{eq:observations}
\end{equation}
Here $\mathcal P_j$ is a source-state probability, not the residual slope
$S_j$, and $\pi_j$ is the selection probability. Overlapping source
attributes require union probabilities rather than double counting.
Even a stationary one-class model constrains a product of emergence
rate, source lifetime $\tau$, and selection probability. Increasing one
factor and decreasing another can preserve the likelihood. A null survey
therefore does not by itself identify an inward transition.

\subsection{A limited thermal bridge}
\label{sec:thermal-bridge}
Landauer's bound assigns at least $k_BT\ln2$ of heat to erasure of one
unbiased bit under its thermodynamic assumptions \citep{Landauer1961}.
It is not a universal energy charge per reversible logical operation or
a law fixing subjective clock speed \citep{Bennett1973}. For an ideal
one-temperature blackbody radiator, negligible background temperature,
total effective emitting area $A=M_{\rm rad}/\Sigma_{\rm rad}$, and
thermalized luminosity $L$, the Stefan--Boltzmann law gives
\begin{equation}
 T=\left(\frac{L\Sigma_{\rm rad}}{\sigma_{\rm SB}M_{\rm rad}}\right)^{1/4}.
 \label{eq:radiator}
\end{equation}
This equation follows by solving $L=\sigma_{\rm SB}AT^4$ for $T$;
$A$ already counts all radiating faces. With nominal solar luminosity,
$\Sigma_{\rm rad}=1\,{\rm kg\,m^{-2}}$ as an assumed engineering
parameter, and the JPL planetary masses \citep{JPLPhysical}, a Mercury
mass gives $378.16\,{\rm K}$ and a Jupiter mass $43.43\,{\rm K}$.
Their wavelength-form Wien peaks are $7.66$ and $66.73\,\mu{\rm m}$,
respectively. Constants and numerical root checks are specified in the
supplement \citep{NIST2022}.

The first example lies inside the $100$--$700\,{\rm K}$ model range
searched by Project Hephaistos II; the second lies below it
\citep{Suazo2024}. The WISE/Spitzer framework of \citet{Wright2014b}
likewise targets mid-infrared waste heat, not all possible cold radiator
architectures. Thus inward allocation does not predict a unique cold
far-infrared population. It permits temperature-dependent search targets
once dissipated power and radiating area are supplied. Bradbury's
Matrioshka-brain manuscript already describes nested stellar computing
and heat disposal \citep{Bradbury1999}. The cooler galactic environments
of \citet{CirkovicBradbury2006} can motivate migration, not necessarily
continued residence around the original star. The present bridge supplies
no priority claim for either idea and no construction-feasibility proof.
